\chapter{Methodology and Evaluation}\label{ch:methods}
This chapter defines the prediction tasks, information flow, and comparisons used to interpret the saved experiments. No models were retrained for the report; numerical checks operate on saved outputs.

\section{Notation and prediction targets}
Let $i$ identify a physical specimen, $t$ an observation time, $I_{it}$ its image, and $m_{it}$ the available design and exposure metadata. Let $v_{it}$ denote a visible-corrosion label, $a_i$ terminal wire-area loss, and $u_i$ terminal ultimate load. A representation $\phi(I_{it})$ converts an image into model inputs.

Visible-corrosion regression estimates
\begin{equation}\widehat v_{it}=f\bigl(\phi(I_{it}),m_{it}\bigr),\end{equation}
with image-only variants omitting metadata. Terminal-capacity regression instead targets
\begin{equation}\widehat u_{it}=g\bigl(\phi(I_{it}),m_{it}\bigr),\qquad y_{it}=u_i.\label{eq:capacitytarget}\end{equation}
The repeated label in Eq.~\eqref{eq:capacitytarget} is terminal load at every training time. The main capacity score is calculated at the measured terminal image of each held-out specimen. It does not measure performance at a fixed early observation horizon.

\begin{table}[htbp]\centering\small
\caption[Information available at prediction time]{Information availability in the focused capacity analysis. The primary evaluation is retrospective at the terminal image. A future early-image experiment must verify that every predictor is available before its chosen cutoff.}\label{tab:availability}
\begin{tabularx}{\textwidth}{p{.27\textwidth}X}\toprule
Quantity & Role and timing boundary\\\midrule
Image descriptors & Derived from the image at its observation time.\\
Design/treatment metadata & Recorded specimen context; availability is an assumption to check for a future deployment.\\
Week and ageing days & Exposure context at image acquisition.\\
Failure-surface cover & Included metadata; pre-test availability is not established by the saved feature definition.\\
Terminal ultimate load & Training target, repeated over image rows; evaluation uses terminal observations. Excluded from predictors.\\
Wire-area loss / visible labels & Excluded predictors in the focused capacity experiment; use in other phases is explicitly distinguished.\\\bottomrule
\end{tabularx}
\end{table}

\section{Three evaluation regimes}\label{sec:regimes}
All train/test partitions must keep a specimen's observations together. The audited saved manifests satisfy this disjointness requirement. Figure~\ref{fig:regimes} illustrates the unit being withheld.

\textbf{Grouped specimen holdout} withholds specimens sampled from the observed experimental mixture. Unstratified grouped splitting does not guarantee matching proportions of campaigns or treatments. The later capacity analysis explicitly stratifies specimen folds by mesh in the pooled case.

\textbf{Leave-one-treatment-out (LOTO)} withholds a treatment group. Its interpretation depends on the treatment definition: the earlier interpretable phase uses seven coarse groups, while the refined surface experiment uses nine more specific groups. Their results therefore cannot be treated as identical tests of a single invariant property.

\textbf{Leave-one-campaign-out (LOCO)} fits on one campaign and evaluates on the other. With two campaigns, there are two directions of transfer. Because campaign, mesh, chloride, and terminal duration coincide, this is a severe extrapolation stress test. The term does not imply a representative sample of all future sites or campaigns.

\fig{evaluation_regimes}{.82}{Evaluation regimes with specimen histories kept intact. Grouped holdout samples specimens from the observed mixture; LOTO withholds a treatment definition; LOCO withholds a campaign and its aligned design/exposure combination. Tile counts are illustrative and do not represent actual fold sizes. The actual manifests, including mesh-stratified and post-onset analyses, were checked for specimen overlap.}{fig:regimes}

\section{Image representations and model families}
The earliest exploratory scripts demonstrate a workflow but supply no held-out benchmark used here. The classical benchmark uses colour statistics and threshold-derived mask descriptors with regression models. The deep representation uses a frozen ResNet-18 encoder and a learned regression head; its two configurations are image-only and image plus metadata. It has no metadata-only baseline.

The interpretable multi-family phase combines spatial rust summaries, colour, texture, and morphology. Its structural models also include metadata and recorded visible labels; they are mixed-input models, not image-only estimators. The later robustness analysis evaluates cleaned feature sets and selects structural configurations using a weighted rank criterion across regimes.

The focused load experiment compares metadata, RGB, HSV, RGB+HSV, and each corresponding metadata combination. Its regressors include Ridge, RandomForest, GradientBoosting, XGBoost, and CatBoost. Categorical metadata are imputed and encoded within training folds; Ridge inputs are standardized using the training data. Image labels, measured structural targets, identifiers, and future-derived summary fields are excluded from the input feature list of this focused analysis. Detailed feature names and source locations appear in Appendix~\ref{app:reproduction}.

\section{Terminal-load comparison design}\label{sec:capacitydesign}
The pooled load experiment uses two repetitions of five specimen folds, covering all 48 specimens in each repetition. Each held-out fold contains nine or ten terminal specimens. Training uses all retained image rows belonging to training specimens, while metrics use the terminal observations of test specimens. The resulting image-row objective weights specimens in proportion to the number of retained photographs unless explicitly corrected; the saved primary workflow does not apply such correction.

Models are selected separately within feature families using the saved evaluation scores. This comparison asks whether image-enhanced configurations consistently improve on a metadata anchor under a matched partition design. Since the same folds help select and report the winners, the results remain exploratory estimates for selected configurations, not a nested evaluation of the complete selection procedure \citep{cawley2010}.

Two sensitivity analyses constrain the pooled interpretation. The first repeats the comparison within each mesh family, leaving 24 independent structural specimens per family and only four or five terminal test specimens per fold. The second retains observations at and after a specimen's first recorded total-rust value of at least 0.5\%, leaving 43 specimens with eligible terminal observations and eight or nine terminal test specimens per fold. The threshold is a pragmatic sensitivity cutoff. The code retains all later observations once onset occurs; it does not require each later value to remain above the cutoff.

A separate specimen-summary analysis first predicts load using metadata and then attempts to correct residuals with corrosion summaries within mesh groups. Its training-stage cross-fitting concerns that refinement experiment. It must not be confused with the in-sample structural estimates that feed the degradation pipeline.

\section{Metrics, aggregation, and uncertainty}\label{sec:metrics}
For a test set of $n$ observations with errors $e_j=\widehat y_j-y_j$,
\begin{equation}
\operatorname{MAE}=\frac{1}{n}\sum_j |e_j|,\qquad
\operatorname{RMSE}=\sqrt{\frac{1}{n}\sum_j e_j^2}.
\end{equation}
MAE measures typical absolute error in the target's units; RMSE places more weight on larger errors. The coefficient of determination is
\begin{equation}
R^2=1-\frac{\sum_j e_j^2}{\sum_j(y_j-\bar y_{\rm test})^2}.
\end{equation}
Negative test $R^2$ means squared error exceeds the test-set-mean reference for that fold. That reference is a score definition, not a fitted deployable baseline. In small, narrow-range subgroups its denominator can be small, making fold $R^2$ unstable. Spearman correlation $\rho$ evaluates rank agreement and supplies complementary information, but is also variable in small folds.

The main load tables report arithmetic means across ten folds and, where shown, sample standard deviations across those folds. Repetitions share specimens and training sets; their standard deviation is not a confidence interval on future error. An average of fold scores also need not equal a score calculated after averaging a specimen's repeated predictions. Neither a small difference nor overlapping descriptive error bars establishes statistical equivalence.

\section{Fixed models versus reselected winners}
Two legitimate summaries answer different questions. A per-regime winner describes the best selected configuration under that regime. A fixed-model robustness comparison follows the same grouped-holdout winner into each stricter regime. Only the latter supports a ratio of regime MAE to that model's own grouped MAE. Cross-generation comparisons in this report use such within-generation ratios or qualitative patterns, rather than a common raw-error ranking.

These definitions keep observations, model selection, and transfer questions separate. The next two chapters apply them first to the visible targets and then to terminal structural outcomes.
